\section{Other possible directions: GPs}

\begin{frame}[more-space]{Other possible directions: GPs}
  \dataheading{Gaussian process regression as a small-data alternative}
  \begin{columns}[T,onlytextwidth]
    \column{0.48\textwidth}
      \begin{beamercolorbox}[sep=2mm,wd=\linewidth]{data panel}
        \dataheading{Why it is worth comparing}
        \footnotesize
        A Gaussian process (GP) uses a kernel to express how nearby operating
        conditions should have similar temperatures.\par
        \smallskip
        It can represent nonlinear response shapes without selecting a fixed polynomial degree.\par
        \smallskip
        It returns a predictive mean and variance conditional on the kernel,
        noise assumptions and training data. A small dataset makes exact GP fitting feasible.
      \end{beamercolorbox}
    \column{0.49\textwidth}
      \begin{beamercolorbox}[sep=2mm,wd=\linewidth]{data warning}
        \dataheading{What still needs validation}
        \footnotesize
        Nineteen points do not guarantee a reliable GP. Kernel and noise choices
        can strongly influence results and extrapolation.\par
        \smallskip
        Predictive uncertainty is model-dependent; it is not automatically a
        calibrated bound on real CPU temperature.\par
        \smallskip
        UQ and Sobol analysis remain possible through sampling/integration;
        they lack the same simple PCE coefficient-group formulas in general.
      \end{beamercolorbox}
  \end{columns}
  \medskip
  {\footnotesize
  A GP could complement the standard PCE and PI-PCE benchmark.
  It does not repair missing input combinations or replace experimental checks.
  }
  \datasource{Primary software reference: \href{https://scikit-learn.org/stable/modules/gaussian_process.html}{scikit-learn Gaussian process regression documentation}.
    GP comparison is an alternative proposed in the notes; no GP has been trained.}
\end{frame}

\begin{frame}[more-space]{Other possible directions: GPs}
  \dataheading{A fair benchmark: same observations, same held-out cases}
  {\footnotesize
  \renewcommand{\arraystretch}{1.3}
  \begin{tabular}{@{}p{0.20\textwidth}p{0.36\textwidth}p{0.38\textwidth}@{}}
    \toprule
    \textbf{Candidate} & \textbf{Main idea} & \textbf{Uncertainty interpretation}\\
    \midrule
    Standard PCE & Data-only polynomial regression, degrees 1/2. & Analytical input-induced moments; no automatic prediction-error model.\\
    Physics-informed PCE & Same basis, with a thermal-resistance loss term. & Same input-UQ formulas; physical guidance is not an error guarantee.\\
    GP regression & Kernel-based nonlinear prediction. & Conditional predictive variance; coverage and noise assumptions need checking.\\
    \bottomrule
  \end{tabular}
  }
  \medskip
  \begin{beamercolorbox}[sep=2mm,wd=\textwidth]{data panel}
    \dataheading{Compare what matters for the final decision}
    \footnotesize
    Leave one experimental case out; tune each candidate inside the other cases.
    Report RMSE, MAE, $R^2$, maximum error and residuals; also compare physical
    plausibility, local underprediction, uncertainty calibration and optimization performance.
    Use equivalent decision scenarios and independently check promising settings.
  \end{beamercolorbox}
  \datasource{Benchmark proposal from the project notes. No ranking, fitted-model results or optimization performance is available yet.
    The earlier 10-fit workload covers only the untuned two-degree standard-PCE baseline; tuning, PI-PCE and GP add fits.}
\end{frame}
